\begin{abstract}
Network-on-Chip (NoC) architectures have become the standard interconnect fabric
for many-core systems, yet most proposals face a fundamental trade-off between
latency, area, and congestion management. This paper presents \textbf{HyNoC}
(Hybrid Network-on-Chip), an open-source NoC architecture that combines
circuit-switch path establishment with wormhole data transfer, targeting
distributed computing systems built around VLIW processor cores on FPGA.

HyNoC employs source routing, where the complete path through the network is
encoded in the packet header by the sender or statically at compile time, enabling
both deterministic low-latency transfers and software-level hotspot avoidance
without the area overhead of virtual channels. The router features a parallel
round-robin arbiter (PRRA) with fixed grant latency, per-port independent clock
domains, and support for both unicast and multicast routing.

We discuss the design rationale, the hybrid switching model, and positioning
relative to prior NoC art, and argue that a richer topology combined with
compiler-assisted static routing is a competitive alternative to virtual channels
for FPGA-based distributed VLIW systems. The complete RTL implementation is
released as open-source hardware under the CERN-OHL-P~v2 license. The Veriparse source-to-source elaboration toolkit~cite{veriparse} is provided as a companion tool to handle advanced RTL constructs and accelerate simulation.
\end{abstract}

\begin{IEEEkeywords}
Network-on-Chip, FPGA, source routing, wormhole switching, circuit switching,
VLIW, parallel round-robin arbiter, open-source hardware.
\end{IEEEkeywords}
